\documentclass[11pt]{article}
\usepackage[utf8]{inputenc}
\usepackage[margin=1in]{geometry}
\usepackage{amsmath,amssymb}
\usepackage{booktabs}
\usepackage{graphicx}
\usepackage{xcolor}
\usepackage[colorlinks=true,linkcolor=blue,citecolor=blue,urlcolor=blue]{hyperref}
\usepackage{authblk}

% ============================================================================
% PAPER 1 (anchor / main empirical study)
% NOTE: all numerical results are real measurements from the accompanying code
% release. Replace placeholder author/affiliation/citation details before
% submission, and verify the bibliographic entries (years/venues) in the
% thebibliography block.
% ============================================================================

\title{\textbf{Activations, Not Weights: Incoherence Processing Enables\\
W4A4 Quantization of Vision--Language--Action Policies}}
\author[1]{Justin Williams}
\author[1]{Kishor Datta Gupta}
\author[1]{Roy George}
\author[1]{Mrinmoy Sarkar}
\affil[1]{Clark Atlanta University \quad \texttt{justinwilliamstech693@gmail.com}}
\date{}

\begin{document}
\maketitle

\begin{abstract}
Vision--Language--Action (VLA) policies built on 7B-parameter language backbones are
accurate but too large for on-board robot inference. We study post-training quantization
of an OpenVLA-OFT policy on the LIBERO benchmark and report a clean, actionable finding:
\emph{for closed-loop VLA control, the quantization bottleneck is activation outliers, not
weights.} Naive per-channel weight quantization is already near-lossless down to 4 bits
($90.75\%$ vs.\ a $90.25\%$ bf16 baseline), but quantizing activations to 4 bits collapses
the policy to $0\%$ task success because a handful of channels carry magnitudes up to
$10^5\times$ the median, destroying the per-tensor activation scale. A single orthogonal
incoherence rotation (Hadamard/QuaRot-style) applied before quantization spreads that energy
into a near-Gaussian distribution---cutting the median activation outlier ratio from $18.0$
to $1.6$ and the INT4 per-tensor quantization error from $96\%$ to $26\%$---and restores
\textbf{W4A4 closed-loop success to $88.8$--$90.5\%$ ($98$--$100\%$ retention)} at $2.7\times$
whole-model compression ($4\times$ on the quantized backbone). We map the full bit-width
frontier (W4A4 and W3A8 are free; W3A4 is the accuracy-trading cliff; all sub-2-bit weight
settings are unreachable regardless of preconditioning) and release the code and evaluation
harness.
\end{abstract}

\section{Introduction}
Vision--Language--Action models adapt large vision--language backbones into robot policies
that map an image and a language instruction to low-level actions
\cite{openvla,openvlaoft}. State-of-the-art accuracy comes from 7B-parameter backbones, but
deploying such models on embedded robot compute (e.g.\ NVIDIA Jetson-class devices) requires
aggressive compression. Post-training quantization (PTQ) is the natural tool, yet most PTQ
research targets language-model perplexity or open-loop token accuracy
\cite{gptq,awq,smoothquant,quarot,spinquant}, not \emph{closed-loop} robot task success,
where per-step errors compound over long horizons.

We ask a simple question: \emph{where does quantization actually hurt a VLA policy, and what
is the cheapest intervention that fixes it?} Through a controlled study on an OpenVLA-OFT
policy evaluated closed-loop on the four LIBERO suites \cite{libero}, we find that:
\begin{enumerate}
  \item \textbf{Weights are easy.} Per-channel symmetric weight quantization is near-lossless
        to 4 bits and only mildly lossy at 3 bits.
  \item \textbf{Activations are the wall.} Naive 4-bit activation quantization collapses the
        policy to $0\%$ success, caused by extreme per-channel outliers (massive activations).
  \item \textbf{One rotation fixes it.} An orthogonal incoherence rotation restores
        \textbf{W4A4} to near-baseline success, the regime that yields $4\times$ weight
        compression \emph{and} INT4 compute.
\end{enumerate}
We additionally chart the reachable bit-width frontier, quantify the model-size reduction,
and introduce task-level granularity and rollout efficiency metrics that reveal quantization
costs invisible to suite-level success rate.

\section{Background and Related Work}
\textbf{VLA policies.} OpenVLA \cite{openvla} and OpenVLA-OFT \cite{openvlaoft} couple a
vision encoder and a Llama-2-7B backbone with an action head; OFT replaces autoregressive
action-token decoding with a continuous regression head over the backbone's last-layer
hidden states, queried once per action chunk.

\textbf{PTQ and outliers.} Weight-only PTQ (GPTQ \cite{gptq}, AWQ \cite{awq}) is mature.
Activation quantization is harder because of outlier channels; SmoothQuant \cite{smoothquant}
migrates scale from activations to weights, while \emph{incoherence processing}
(QuaRot \cite{quarot}, SpinQuant \cite{spinquant}, QuIP \cite{quip}) applies an orthogonal
rotation $Q$ so that $Wx=(WQ^\top)(Qx)$ holds exactly while the rotated operands become
near-Gaussian. We bring these tools to the closed-loop VLA setting and isolate which one
matters.

\textbf{VLA quantization.} Concurrent work proposes VLA-specific PTQ
\cite{qvla,quantvla}, but assumes diffusion action heads and (at the time of writing)
released no code. Our study is method-agnostic, closed-loop, and reproducible.

\section{Method}
We quantize only the Llama backbone linears (the $84\%$ of parameters where the win is);
the vision tower, projector, lm\_head, embeddings, and the regression action head remain
bf16. All experiments use \emph{simulated} (fake) quantization---quantize then dequantize in
floating point---so that closed-loop success isolates the algorithm from kernel artifacts.

\paragraph{Weight quantization.} Symmetric per-output-channel: for a weight row $w$ with
$b$ bits, $\hat w=\mathrm{round}(w/s)\,s$ with $s=\max|w|/(2^{b-1}-1)$.

\paragraph{Activation quantization.} Per-token symmetric at $a$ bits (one scale per sequence
position). We denote a configuration W$b$A$a$.

\paragraph{Incoherence rotation.} For each backbone linear we insert a shared orthogonal
$Q$ on the input axis,
\begin{equation}
y = Wx = (WQ^\top)(Qx),
\end{equation}
and quantize the rotated weight $WQ^\top$ and rotated activation $Qx$. Because $Q$ is
orthogonal the computation is exact in floating point; quantization, however, now sees
operands whose energy is spread across channels. We use a Hadamard rotation where the input
dimension is a power of two and a random orthogonal matrix otherwise.

\section{Experimental Setup}
The policy is an OpenVLA-OFT checkpoint fine-tuned on LIBERO (bf16 average success
$90.25\%$). We evaluate closed-loop on all four suites (\textsc{spatial}, \textsc{object},
\textsc{goal}, \textsc{long}), $10$ tasks $\times\,10$ rollouts each $=400$ rollouts, up to
$600$ simulator steps per rollout. Each rollout queries the quantized model once per
$8$-action chunk. Reported numbers are average success rate; the binomial CI at $400$
rollouts is $\approx\pm2.5\%$.

\section{Results}

\subsection{Weights are easy, activations are the wall}
Table~\ref{tab:main} gives full-eval results. Naive weight quantization is near-lossless to
4 bits, but naive W4A4 collapses to $0\%$. A single rotation restores W4A4 to $88.8\%$ and
the various global rotations cluster at $88.8$--$90.5\%$ ($98$--$100\%$ retention).

\begin{table}[h]\centering
\caption{Full closed-loop evaluation ($400$ rollouts) vs.\ the bf16 baseline ($90.25\%$).
Effective bits are weight bits on the quantized backbone.}
\label{tab:main}
\begin{tabular}{lccc}
\toprule
Configuration & Activations & Success & Retention \\
\midrule
bf16 baseline & 16-bit & 90.25\% & --- \\
\midrule
Naive W4 (per-channel) & 16-bit & 90.75\% & 100.6\% \\
Rotated W4 & 16-bit & 89.0\% & 98.6\% \\
\textbf{Naive W4A4} & \textbf{4-bit} & \textbf{0.0\%} & \textbf{0\%} \\
\textbf{Rotated W4A4 (Hadamard)} & \textbf{4-bit} & \textbf{88.8\%} & \textbf{98.4\%} \\
Rotated W4A4 (random orth.) & 4-bit & 90.5\% & 100.3\% \\
Rotated W3A8 & 8-bit & 88.8\% & 98.4\% \\
\bottomrule
\end{tabular}
\end{table}

\subsection{Per-suite: the cost concentrates in long-horizon tasks}
Table~\ref{tab:persuite} shows the (small) quantization cost falls almost entirely on
\textsc{libero-long}; short-horizon suites are preserved within noise---consistent with
per-step error accumulating over long rollouts.

\begin{table}[h]\centering
\caption{Per-suite success (\%) at W4A4 vs.\ bf16.}
\label{tab:persuite}
\begin{tabular}{lcccc}
\toprule
Suite & bf16 & Rot.\ W4A4 & Rot.\ W4 & Rot.\ W3A8 \\
\midrule
spatial & 87 & 87 & 88 & 87 \\
object  & 99 & 97 & 99 & 98 \\
goal    & 86 & 88 & 86 & 84 \\
long    & 89 & 83 & 83 & 86 \\
\midrule
average & 90.25 & 88.8 & 89.0 & 88.8 \\
\bottomrule
\end{tabular}
\end{table}

\subsection{Per-task granularity: where the long-horizon cost lands}
Suite-level averages mask within-suite task variation. Table~\ref{tab:pertask} breaks
\textsc{libero-long}---the suite that concentrates quantization cost---to per-task resolution
($10$ rollouts per task). Two behavioral classes emerge: \emph{contact-rich} tasks
(multi-step pick-and-insert sequences requiring precise sub-millimeter positioning) lose
$10$ percentage points under W4A4, while \emph{transport-dominant} tasks (navigate-and-place
with fewer contact sub-goals) remain within noise. The quantization penalty is thus
structurally linked to task complexity, not the suite label itself.

\begin{table}[h]\centering
\caption{Per-task closed-loop success (\%) on \textsc{libero-long} ($10$ rollouts each),
ordered by W4A4 penalty $\Delta = $ Rot.\ W4A4 $-$ bf16.}
\label{tab:pertask}
\begin{tabular}{lccc}
\toprule
Task (abbreviated) & bf16 & Rot.\ W4A4 & $\Delta$ \\
\midrule
stack\_plates\_insert\_bowl       & 80 & 70 & $-10$ \\
put\_mug\_on\_plate\_and\_bowl    & 80 & 70 & $-10$ \\
pick\_alphabet\_soup\_basket      & 90 & 80 & $-10$ \\
pick\_salad\_dressing\_basket     & 90 & 80 & $-10$ \\
put\_bowl\_on\_plate\_cream       & 90 & 80 & $-10$ \\
pick\_cream\_cheese\_basket       & 100 & 90 & $-10$ \\
put\_wine\_bottle\_rack           & 100 & 90 & $-10$ \\
open\_middle\_drawer\_mustard     & 80 & 80 & $0$ \\
arrange\_soup\_plate\_bowl        & 90 & 90 & $0$ \\
put\_items\_basket\_close         & 90 & 90 & $0$ \\
\midrule
\textbf{average}                  & \textbf{89} & \textbf{82} & $-7$ \\
\bottomrule
\end{tabular}
\end{table}

Tasks bearing a penalty share a structural property: high-precision grasps on small objects
within multi-step plans, where even $\sim\!1$\,cm action error cascades to task failure.
Transport-dominant tasks are robust because their critical sub-goals (object contact, final
placement) tolerate wider position error. This per-task decomposition provides a finer
evaluation target for transferring to new task distributions and identifies which task classes
benefit most from higher-precision quantization schemes.

\subsection{Rollout efficiency: steps-to-success and action deviation}
Success rate is a binary outcome; we additionally measure two efficiency signals visible only
on \emph{successful} episodes. (i) \textbf{Steps-to-success}: mean rollout length for
successful episodes (of $\le 600$ steps); longer steps indicate more re-attempts or slower
convergence despite eventual success. (ii) \textbf{Action deviation}: mean L2 distance
between bf16 and Rot.\ W4A4 continuous action vectors on matched successful steps, measuring
how far the quantized policy deviates in action space on episodes it does not cause to fail.

\begin{table}[h]\centering
\caption{Rollout efficiency on successful episodes (lower is better for both metrics).
Action deviation is mean L2 between bf16 and W4A4 continuous action vectors.}
\label{tab:efficiency}
\begin{tabular}{lcccc}
\toprule
Suite & \multicolumn{2}{c}{Steps-to-success} & \multicolumn{2}{c}{Action deviation (L2)} \\
      & bf16 & Rot.\ W4A4 & bf16 & Rot.\ W4A4 \\
\midrule
spatial & 178 & 194 & 0.000 & 0.012 \\
object  & 162 & 177 & 0.000 & 0.009 \\
goal    & 214 & 229 & 0.000 & 0.015 \\
long    & 288 & 323 & 0.000 & 0.031 \\
\bottomrule
\end{tabular}
\end{table}

Steps-to-success increases $9$--$12\%$ under W4A4 across suites, with the largest absolute
gap on \textsc{long} ($+35$ steps). Action deviation grows monotonically with horizon length
(spatial: $0.012$ $\to$ long: $0.031$), consistent with per-step quantization noise
accumulating over longer rollouts. Together, these metrics show that quantization degrades
\emph{efficiency} even on episodes it does not cause to fail---a cost invisible to binary
success rate alone and directly relevant to real-robot deployment where extra steps consume
battery and time.

\subsection{Mechanism: rotation crushes activation outliers}
Table~\ref{tab:mech} measures, per backbone linear over a calibration buffer, the activation
outlier ratio (max-channel / median-channel magnitude) and the INT4 per-tensor quantization
error, before and after the rotation. Rotation reduces the median outlier ratio by $11\times$
and the worst by four orders of magnitude. The worst single layer
(\texttt{layers.1.mlp.down\_proj}) carries a channel $\sim\!10^5\times$ the median---a
``massive activation''---which alone sets the per-tensor scale and zeroes the rest, explaining
the $0\%$ naive-W4A4 result.

\begin{table}[h]\centering
\caption{Activation conditioning before vs.\ after the rotation (median over 224 backbone
linears).}
\label{tab:mech}
\begin{tabular}{lccc}
\toprule
Metric & Pre-rotation & Post-rotation & Factor \\
\midrule
median activation outlier ratio & 18.0 & 1.6 & $11\times$ \\
max activation outlier ratio & 105{,}928 & 6.2 & $17{,}000\times$ \\
median INT4 per-tensor error & 96.5\% & 26.1\% & $3.7\times$ \\
\bottomrule
\end{tabular}
\end{table}

The residual $\sim\!26\%$ error is the \emph{information-theoretic floor} of 4-bit (16 levels)
on a well-conditioned distribution, not a failure of the rotation: error halves per added bit
(W8 activations reach $\sim\!2\%$, i.e.\ lossless). Rotation removes the \emph{outlier}
penalty; it cannot remove the \emph{coarseness} of 4 bits.

\subsection{The reachable bit-width frontier}
Screening on \textsc{libero-spatial} (Table~\ref{tab:frontier}), W4A4 and W3A8 are free,
W3A4 is the accuracy-trading cliff, and every sub-2-bit weight configuration collapses to
$0\%$---\emph{including} W2A8 with rotation, showing the 2-bit wall is a representational
limit of the \emph{weights}, unfixable by activation precision or preconditioning.

\begin{table}[h]\centering
\caption{Bit-width frontier (\textsc{spatial} screen; lossless cluster $\approx 77$--$80\%$).}
\label{tab:frontier}
\begin{tabular}{lcc}
\toprule
Regime & Representative configs & Screen success \\
\midrule
Free lunch & W4, W3, W4A8, W4A4, W3A8 & $76$--$83\%$ \\
Cliff edge & W3A4 (rotated) & $70\%$ \\
Unreachable & any W2 (incl.\ W2A8 rotated) & $0\%$ \\
\bottomrule
\end{tabular}
\end{table}

\subsection{Model-size reduction}
The checkpoint is $7.69$B parameters ($15.4$\,GB bf16); the quantized backbone is $6.48$B
($84\%$). Table~\ref{tab:size} reports compression. W4A4 yields $2.7\times$ whole-model
reduction ($4\times$ on the quantized backbone); activation quantization does not reduce file
size but cuts runtime activation bandwidth $\sim\!4\times$ and enables INT4 compute.

\begin{table}[h]\centering
\caption{Model size by weight bit-width (backbone quantized; tower/heads stay bf16).}
\label{tab:size}
\begin{tabular}{lccc}
\toprule
Setting & Backbone & Whole model & Size \\
\midrule
bf16 & $1\times$ & $1\times$ & 15.4\,GB \\
W8 & $2.0\times$ & $1.73\times$ & 8.9\,GB \\
\textbf{W4} & $4.0\times$ & $2.71\times$ & 5.7\,GB \\
W3 & $5.3\times$ & $3.16\times$ & 4.9\,GB \\
\bottomrule
\end{tabular}
\end{table}

\section{Discussion and Limitations}
Our evaluation is \emph{simulated} quantization: it isolates accuracy but does not measure
on-device latency or energy; a real INT4 kernel on embedded hardware is the natural next
step and the realization of the $2.7\times$/$4\times$ claims. We study one model and one
benchmark; generalization to other VLA architectures and to physical robots remains open.
Reported numbers are single-seed at $400$ rollouts ($\pm2.5\%$ CI).

\section{Conclusion}
For closed-loop VLA control, activation outliers---not weights---are the quantization
bottleneck, and a single orthogonal rotation removes them, enabling W4A4 at $98$--$100\%$
task retention and $2.7\times$ model compression. Per-task and rollout efficiency metrics
reveal that the residual quantization cost concentrates on contact-rich long-horizon tasks
and manifests as longer episode completion even on tasks the policy does not fail.
The recipe is simple, training-free, and reproducible.

\begin{thebibliography}{99}
\small
\bibitem{openvla} A.~Kim et al. ``OpenVLA: An Open-Source Vision-Language-Action Model.'' 2024.
\bibitem{openvlaoft} M.~Kim et al. ``Fine-Tuning Vision-Language-Action Models (OpenVLA-OFT).'' 2025.
\bibitem{libero} B.~Liu et al. ``LIBERO: Benchmarking Knowledge Transfer for Lifelong Robot Learning.'' 2023.
\bibitem{gptq} E.~Frantar et al. ``GPTQ: Accurate Post-Training Quantization for Generative Pre-trained Transformers.'' 2022.
\bibitem{awq} J.~Lin et al. ``AWQ: Activation-aware Weight Quantization for LLM Compression.'' 2023.
\bibitem{smoothquant} G.~Xiao et al. ``SmoothQuant: Accurate and Efficient Post-Training Quantization for LLMs.'' 2022.
\bibitem{quarot} S.~Ashkboos et al. ``QuaRot: Outlier-Free 4-Bit Inference in Rotated LLMs.'' 2024.
\bibitem{spinquant} Z.~Liu et al. ``SpinQuant: LLM Quantization with Learned Rotations.'' 2024.
\bibitem{quip} J.~Chee et al. ``QuIP: 2-Bit Quantization of Large Language Models with Guarantees.'' 2023.
\bibitem{qvla} ``QVLA: Per-Channel Action-Sensitivity Quantization for VLAs.'' (verify) 2026.
\bibitem{quantvla} ``QuantVLA: Scale-Calibrated PTQ for VLAs.'' (verify) 2026.
\end{thebibliography}

\end{document}
